\section{Related Work}

Vibe Theory sits inside a long and serious tradition. For decades careful researchers have asked whether
spacetime is fundamentally discrete, whether physics is computation, and whether geometry comes from
information. Vibe Theory does not invent that question. It joins a conversation that already includes causal
sets, causal dynamical triangulations, loop quantum gravity, the Wolfram Physics Project, the
cellular-automaton interpretation of quantum mechanics, quantum cellular automata, holographic
error-correcting codes, hyperbolic lattices, hyperbolic cellular automata, and the informational
reconstructions of quantum mechanics.

Much of what Vibe Theory needs has already been proven by other
people. The framing we adopt is that Vibe Theory assembles established results onto one specific substrate
and then asks one further question those programs mostly leave aside. What is the felt interior of the thing,
and what direction does it move in.

This section places Vibe Theory among those programs. For each family it states what is shared, what differs,
and what Vibe Theory adds, and it points to the specific experiment that touches each one. We try to credit
the prior work generously and to claim for Vibe Theory only what it earns. In particular, the quantum-field
route that Vibe Theory uses is not a new discovery. It is the established quantum-cellular-automaton and
quantum-walk derivation of the Dirac equation, applied on a hyperbolic substrate. All experiments referenced
here are recorded in the accompanying codebase \cite{pollard2026vibetest}.

\subsection{Causal sets}

Causal-set theory, in the line of Bombelli, Lee, Meyer, and Sorkin \cite{bombelli1987} and the classical
sequential-growth dynamics of Rideout and Sorkin \cite{rideout1999}, takes the most radical version of the
discrete idea seriously. Spacetime is a discrete set of events with a partial order, and order plus number is
all there is. Surya's review collects the program's reach and its open problems \cite{surya2019}. Vibe Theory
shares the core commitment. Space and time are read off a discrete relational structure, and distance is
step-counting on a graph rather than a length measured by an external ruler.

The difference is the price of Lorentz invariance. Causal-set theory pays for it with randomness. A Poisson
sprinkling of Minkowski space has no preferred direction in distribution, so the discreteness does not pick
out a rest frame. The cost is that the substrate is not a lawful, repeatable structure. Vibe Theory takes the
opposite route. It keeps an exact, ordered, deterministic crystal and gets Lorentz safety from curvature
instead of from disorder. A flat lattice has a grain that announces a rest frame. A negatively curved crystal
scrambles directions under transport, so it looks frameless from inside even though every cell is identical.

What Vibe Theory adds is the demonstration that one need not choose between lawfulness and Lorentz invariance.
Experiment P84 reproduces the causal-set result itself. A Poisson sprinkle gives a flat rapidity distribution
that stays flat under a genuine boost, while the matched flat lattice gives a peaked one. Then P40 and P45
show every hyperbolic substrate tested, including the three-dimensional $\{5,3,4\}$ dodecagrid, lands in the
same low-anisotropy band as the random sprinkle, while the flat lattice fails.

So Vibe Theory keeps the
determinism causal sets give up and still passes the test that motivated their randomness. Curvature rather
than disorder buys Lorentz invariance, and that is the bridge between the two programs.

\subsection{Causal dynamical triangulations}

Causal dynamical triangulations, in the work of Ambjorn, Jurkiewicz, and Loll \cite{ambjorn2004}, build
spacetime from small simplices glued along a fixed causal slicing and sum over the geometries. The central
result is that a sensible four-dimensional spacetime emerges from the dynamics rather than being put in by
hand. Loll's review gives the current picture \cite{loll2019}. Vibe Theory shares the goal of letting
dimension and large-scale geometry be outputs of a discrete dynamics rather than postulates, and the
insistence on a causal, time-ordered substrate rather than a Euclidean one.

The difference is method. Causal dynamical triangulations is a path integral. It sums over many triangulated
geometries weighted by an action, so the geometry is an ensemble average. Vibe Theory does the opposite. It
fixes one rigid geometry, the $\{5,3,4\}$ crystal, and never sums over alternatives. The dynamics lives in
the tone on that fixed lattice, not in a fluctuating triangulation, and Vibe Theory commits to negative
curvature where the triangulations program lets effective curvature emerge from the averaged ensemble.

What Vibe Theory adds is a single committed geometry on which spectral and dimensional quantities are measured
directly rather than averaged over a path integral. The two programs make opposite methodological bets.
Triangulations trust the ensemble to find the geometry. Vibe Theory trusts a forced geometry and measures
on it.

\subsection{Loop quantum gravity and spin networks}

Loop quantum gravity, and the spin networks of Rovelli and Smolin \cite{rovelli1995}, quantize geometry
itself. Area and volume become operators with discrete spectra, and the quantum state of space is a graph
with spins on its edges. Geometry is built from a combinatorial, graph-theoretic object, and continuous
space is a coarse-grained description of it. Vibe Theory shares the picture that space is a graph of relations
with discrete capacity per cell, and that the smooth manifold is emergent.

The difference is the status of the graph. In loop quantum gravity the spin network is the quantum state of
geometry, and the graph and its labels fluctuate under the dynamics. Vibe Theory holds the graph fixed and
puts the only dynamical variable, the ternary tone, on the cells. So where spin networks make geometry the
quantum degree of freedom, Vibe Theory makes geometry the fixed stage and the tone the actor on it.

What Vibe Theory adds is a fixed graph on which the quantum field is a separate, measured object rather than
identical to the geometry. This keeps the substrate analysis and the field analysis cleanly apart. The crystal
is studied as pure structure, and the quantum behavior is studied as the tone's dynamics on it (P151, P158).

A single thread runs through this whole family. Each program has to answer how a discrete structure can look
the same to every observer. Causal sets answer with random sprinkling, triangulations with an averaged
ensemble, loop quantum gravity with a fluctuating quantum geometry, and Vibe Theory with hyperbolic
regularity. The negative curvature of the $\{5,3,4\}$ crystal scrambles transported directions so that no rest
frame is visible from inside, even though the lattice is exact and deterministic (P40, P45, P84). Vibe Theory
does not claim this is the only way to buy Lorentz safety. It claims it is a way that keeps the substrate exact
and repeatable, which the random and ensemble routes give up.

\subsection{The Wolfram Physics Project}

The Wolfram Physics Project \cite{wolfram2020}, building on the earlier program of \cite{wolfram2002}, is the
closest neighbor in spirit. It builds the universe from a graph or hypergraph and a local rewriting rule,
treats the rule's repeated application as time, reads space off the connectivity rather than positing it, and
reads causal structure off the dependency graph of rewrite events. Vibe Theory agrees on every word of that.
The crystal is a graph, cells are vertices, shared faces are edges, the universe is one local rule running on
it, and the run is time. Both programs also insist that the drawing is not the content. Only the connection
pattern is real.

The deepest difference is the status of the graph. In the Wolfram program the hypergraph grows and rewrites
itself. Nodes and edges are created and destroyed by the rule, the graph's structure is the dynamical
variable, and a vast space of candidate rules is searched for one whose emergent behavior matches physics.
Vibe Theory fixes the graph. The $\{5,3,4\}$ dodecagrid is not searched for and not rewritten. It is held
constant, and the only dynamical variable is a ternary tone living on each fixed cell, evolving by one
conserving rule.

Vibe Theory argues the substrate is forced rather than chosen, by curvature, dimension, and a
minimality preference, where the Wolfram program embraces the multiplicity of rules and studies the whole
rulial space. Wolfram's updating-order ambiguity, the multiway system, becomes a central object for him. Vibe
Theory sidesteps it by putting the dynamics in the tone on a rigid lattice generated once and for all by
reflections rather than discovered by rewriting search.

What Vibe Theory adds is a committed, frozen substrate with a stable local alphabet, which makes the emergent
physics measurable on the exact structure rather than dependent on which rewrite rule was selected. The two
programs are best read as complementary bets. Wolfram bets the right structure is found by searching a vast
rule space. Vibe Theory bets it is forced by a few demands and that the interesting freedom lives in the tone
and the arrow, not in the graph.

\subsection{The cellular-automaton interpretation of quantum mechanics}

The cellular-automaton interpretation of quantum mechanics \cite{thooft2014} argues that beneath the quantum
formalism there is a deterministic, reversible classical automaton, and that quantum states are statistical
descriptions of its underlying ontological states. The program takes seriously that a discrete deterministic
rule can carry the full content of quantum theory once the right basis and the right unitary are identified.
Vibe Theory is squarely in this spirit. Its perception rule, read strictly, is a deterministic reversible
block automaton, and its quantum character appears only when that rule is given its unitary completion (P151,
P126). The conviction that quantum mechanics can sit on top of a deterministic discrete substrate is t
Hooft's, and Vibe Theory inherits it rather than originating it.

The difference is specificity of realization. The cellular-automaton interpretation works largely at the level
of principle and toy models, mapping classes of deterministic systems to quantum ones. Vibe Theory commits to
one concrete substrate and one concrete rule, then measures the resulting dispersion and interference. What it
adds is a specific deterministic-reversible rule on a forced curved lattice whose unitary completion is
measured to give Dirac dispersion and the Born rule, with the reversibility preconditions checked directly
(P126, P128, P131). This is a concrete instance of the thesis on a substrate chosen for independent geometric
reasons, not a generic toy.

\subsection{Quantum cellular automata and quantum walks to the Dirac equation}

This is the family Vibe Theory leans on most heavily, and it is the deepest overlap, so it is important to be
exact about the debt. A line of rigorous work shows that a discrete, local, unitary evolution can have free
quantum field theory as its continuum limit.

\begin{itemize}
\item D'Ariano and Perinotti derive free quantum field theory and the
Dirac dynamics from a denumerable causal network of quantum systems under locality, homogeneity, and
unitarity, with no spacetime assumed in advance \cite{dariano2014}.
\item Bisio, D'Ariano, and Tosini derive the
Dirac equation as the small-wavevector limit of a quantum cellular automaton \cite{bisio2015}, and the program
extends to the Weyl, Dirac, and Maxwell fields in higher dimensions \cite{bisio2016}.
\item Arrighi, Nesme, and Werner establish the unitarity, causality, and
localizability of quantum cellular automata, the theorem that a local unitary evolution is exactly a
finite-depth circuit of local gates \cite{arrighi2011}, and Arrighi, Nesme, and Forets establish the discrete
Lorentz covariance of quantum walks \cite{arrighi2014}.
\item Farrelly and Short show causal fermions in discrete
spacetime arise from the same quantum-walk construction \cite{farrelly2014}.
\end{itemize}

The Arrighi program supplies the structural backbone. The foundational fact, going back
to the coined quantum walk, is that the unitary version of nearest-neighbor spin exchange spreads ballistically
and its continuum limit is the Dirac equation, with the coin angle setting the mass.

The difference is mostly the substrate. This literature works mostly on regular flat lattices or abstract
Cayley graphs. Vibe Theory runs the same construction on the hyperbolic $\{5,3,4\}$ dodecagrid, where the
curvature additionally delivers Lorentz safety, which a flat lattice cannot. A second difference is that the
program generally assumes unitarity, homogeneity, and locality as axioms of the network, where Vibe Theory
derives the unitary object as the completion of a tone rule whose classical, decohered version it also
exhibits.

What Vibe Theory adds must be stated carefully. Vibe Theory does not discover the quantum-walk-to-Dirac route.
That result is established and belongs to the authors above. P151 is the established route applied on the
hyperbolic substrate, not a new derivation of it. What Vibe Theory contributes is to show that its own
perception rule, read strictly, is exactly this construction, and to run it on the forced crystal.

The
stochastic version of the rule gives a massive diffusive field (P136, P137), the classical decohered shadow, a
feature the quantum-cellular-automaton program does not usually exhibit alongside the quantum one. Removing the
illegitimate coin flip leaves a deterministic reversible charge-conserving block automaton that propagates
ballistically (P148, P149). Its unitary completion is a coined quantum walk, and P151 measures the Dirac
dispersion directly, with a ballistic exponent near one, a tunable mass from the coin angle, and reflection
positivity by construction. P154 confirms the exact massless lightcone, and P158 the interference and the Born
rule, with the reversibility preconditions measured directly (P126, P128, P131).

So the contribution is the
hyperbolic substrate, the classical-stochastic shadow, and the forced geometry, plus a demonstration that the
established route is identical to Vibe Theory's stripped rule.

A focused subfamily studies quantum walks on Cayley graphs, whose vertices are group elements and whose edges
are generators. Lopez Acevedo and collaborators classify scalar and coined quantum walks on Cayley graphs and
connect their continuum limits to relativistic equations \cite{lopezacevedo2006}. This matters because the
$\{5,3,4\}$ crystal is generated by a Coxeter group of reflections, which makes it a Cayley-like graph of that
group. The cells are group elements, the moves to neighbors are generators, and a walk on the crystal is a
walk on a Cayley graph in exactly the sense this literature studies.

That literature is usually agnostic about
which group and often works with flat or free groups. Vibe Theory's group is the specific hyperbolic
reflection group of the dodecagrid, which carries the negative curvature and exponential growth that flat or
abelian Cayley graphs lack. The walk theory is borrowed. The hyperbolic Cayley arena is the contribution.

\subsection{Holography, tensor networks, and holographic codes}

Holography is the idea that the information in a region is encoded on its boundary, that a higher-dimensional
gravitational bulk is dual to a lower-dimensional theory on the edge. Maldacena's AdS/CFT correspondence
\cite{maldacena1998}, the Ryu-Takayanagi formula for entanglement entropy as a bulk minimal area
\cite{ryu2006}, Van Raamsdonk's argument that spacetime connectivity is built from entanglement
\cite{vanraamsdonk2010}, and Swingle's identification of the MERA tensor network with a discretized
hyperbolic geometry \cite{swingle2012} together make the case that geometry comes from entanglement and that
the natural home for it is a negatively curved bulk.

A discrete, code-theoretic wing of this program is even
closer to Vibe Theory. The HaPPY pentagon code of Pastawski, Yoshida, Harlow, and Preskill realizes
holography as a quantum error-correcting code on a hyperbolic tiling, where a boundary region reconstructs a
bulk operator \cite{pastawski2015}. Jahn and Eisert review the whole holographic tensor-network field
\cite{jahn2021}. Vibe Theory shares the geometric setting completely. Its substrate is hyperbolic, its bulk
is the anti-de-Sitter-like interior of the $\{5,3,4\}$ crystal, and it treats the boundary-bulk relation as
physical.

The difference is that the holographic codes and tensor networks above are largely static. A fixed network of
tensors on a fixed hyperbolic tiling realizes a fixed code, and the analysis is about the equilibrium
entanglement structure of that state. Vibe Theory starts from a fixed combinatorial crystal but lets a rule
run on it, and asks whether the area law, the bulk shortcut, and a working code appear as measured properties
of the running dynamics rather than as the design of a chosen static network. So Vibe Theory's holography is a
dynamical cousin of the static holographic codes, on a hyperbolic bulk built by reflection rather than wired
by hand.

What Vibe Theory adds is direct measurement on the exact crystal, and a dynamical, growing version of the
holographic code. Experiment P91 measures an entanglement-entropy-like quantity for boundary regions on the
hyperbolic crystal and recovers the Ryu-Takayanagi log law, the entropy set by the length of the subtending
bulk geodesic, with a fit quality near one. The geodesic is a genuine shortcut, and depth into the bulk plays
the role of scale, exactly as in the MERA picture.

Experiments P105 and P127 treat the crystal as a
holographic memory and a growing holographic code, where the code is accreted as the structure grows, the
dynamical counterpart of the static HaPPY construction. P111 then measures the tiny $O(\log N)$ diameter. On a
hundred thousand exact cells no two are more than about six hops apart through the bulk. The same shortcut that
lets entropy ride a geodesic lets a far region of the emergent surface be reached almost instantly through the
interior.

\subsection{Hyperbolic lattices in physics}

A fast-growing experimental and theoretical area puts real physics on hyperbolic lattices. Kollar,
Fitzpatrick, and Houck realize hyperbolic lattices in circuit quantum electrodynamics, building negatively
curved tilings on a superconducting chip \cite{kollar2019}. Maciejko and Rayan develop hyperbolic band
theory, the analog of Bloch theory for these lattices \cite{maciejko2021}. Together these show the hyperbolic
lattice is not a mathematical curiosity but a place where real and simulated physics is being done. Vibe
Theory's substrate is exactly such a lattice, the $\{5,3,4\}$ dodecagrid.

The difference is scope. This literature mostly studies specific phenomena, band structure, spectra, or scalar
fields, on a hyperbolic lattice taken as a given backdrop. Vibe Theory proposes the hyperbolic lattice as the
fundamental substrate of everything, not as a platform for one phenomenon, and runs a single conserving rule on
it. What it adds is the claim that the lattice is fundamental rather than a backdrop, together with a grounding
connection to real experiment. The circuit realizations mean the geometry Vibe Theory bets on is one that
laboratories can build and probe, and its own measurements of dispersion, entropy, and diameter on the
$\{5,3,4\}$ crystal sit alongside the band-theory results as different questions asked of the same geometry.

\subsection{Hyperbolic cellular automata}

Margenstern's body of work proves that cellular automata on hyperbolic tilings are computationally universal,
including a ternary universal automaton on the heptagrid \cite{margenstern2008} and universality across
hyperbolic spaces generally using the railway model \cite{margenstern2013}, where a locomotive runs forever
along tracks and reading-and-flipping switches is the computation, because gliders cannot be choreographed
reliably in hyperbolic space the way they can on a flat grid. Margenstern's splitting method also supplies
the Fibonacci addressing Vibe Theory uses for navigation. Vibe Theory shares the substrate exactly. Its mesh
is a hyperbolic tiling in Margenstern's family, and the $\{5,3,4\}$ dodecagrid is one of the tilings these
theorems cover.

The difference is that Margenstern's universality is a theorem about an abstract automaton wired onto the
tiling, separate from any particular physical dynamics. Vibe Theory needs universality on its own perception
rule, the rule that also produces the quantum field. What it adds is two confirmations of universality on the
live dynamics, anchored to Margenstern's theorem.

\begin{itemize}
\item Experiment P44 shows functional completeness on the running
rule. The ternary signed-majority update is exactly NAND at the right settings, and gates built as real
subgraphs settle to the correct answer as stable fixed points, including a six-gate XOR.
\item Experiment P141 runs
a full Turing-universal Minsky machine with registers stored as charge in the mesh, computing multiplication
using only create, annihilate, and count, the substrate's own primitives.
\end{itemize}

So Margenstern supplies the theorem
on the exact tiling, and P44 and P141 show that Vibe Theory's own rule realizes that universality. Vibe Theory
is careful that universality settles structure only, never feeling, which it treats separately.

\subsection{Hyperbolic graphs and geometric group theory}

A mathematical literature studies the graph-theoretic and large-scale geometry of hyperbolic spaces. Cannon's
exposition of hyperbolic geometry \cite{cannon1997} and Hamann's results on tree-likeness in hyperbolic
graphs \cite{hamann2011} establish that negatively curved graphs are expanders, are tree-like at large scale,
and have exponential ball growth. Vibe Theory uses these facts as the backbone of its substrate analysis. The
crystal graph is a twelve-regular expander, it is tree-like, and its balls grow exponentially at the golden
ratio, which together force a tiny diameter and underwrite the Fibonacci addressing.

The difference is that this literature is pure geometry and combinatorics. It describes the structure without
putting dynamics or physics on it. Vibe Theory turns the same structural facts into physical consequences, a
speed limit from degree-twelve locality, an emergent metric, fast integration, and a holographic readout, and
confirms them by measurement on the exact crystal.

P49 shows the $\{5,3,4\}$ reads like a featureless foam
locally yet is tree-like in its connectivity, unlike a flat lattice. P103 shows the dynamics survive to a
million cells on a twelve-regular hyperbolic expander, because expansion holds at every scale. P111 and P139
measure the logarithmic diameter and the fast $O(\log N)$ coordination the geometry predicts. So Vibe Theory
imports the established hyperbolic-graph mathematics and adds the dynamical measurements and the physical
reading.

\subsection{Informational reconstructions of quantum mechanics}

Chiribella, D'Ariano, and Perinotti derive the full mathematical framework of quantum theory from a short
list of information-theoretic principles, most notably purification, the requirement that every mixed state
arises from a pure state of a larger system \cite{chiribella2011}. This work shows that quantum theory is not
an arbitrary formalism but follows from reasonable demands on how information behaves. Vibe Theory shares the
conviction that quantum structure should be derived, not postulated, and that the right starting point is a
structural or informational one rather than a Hilbert space assumed from the outset.

The difference is that the informational reconstructions are axiomatic and substrate-free. They tell you what
any theory satisfying the principles must look like, without committing to a physical realization. Vibe Theory
is the opposite. It commits to a concrete substrate and rule, lets the quantum structure fall out of the
dynamics, then checks whether the resulting object is genuinely quantum.

What it adds is a specific physical
model whose quantum character is measured rather than assumed. It exhibits a running rule whose unitary
completion shows the two signatures the reconstructions take as targets, interference and the Born rule (P158),
with complex amplitudes that destructively cancel where no classical process can, while the total norm stays
one so that the squared amplitude is a genuine probability. So where Chiribella and colleagues say what
principles force quantum theory, Vibe Theory shows one concrete substrate on which those features appear as
outputs. The two are complementary, the axiomatic why and the constructive how.

\subsection{Consciousness theory}

Every program above is, by intent, a theory of structure. None of them addresses why any arrangement of
parts should be felt from the inside. Chalmers named this the hard problem of consciousness, the gap that
remains even after a complete structural and functional account is given \cite{chalmers1995}. Integrated
information theory, developed by Tononi, attempts to make experience precise by identifying consciousness with
integrated information, a measurable quantity capturing how much a system's whole exceeds the sum of its parts
in a way that cannot be decomposed \cite{tononi2004}. The older informational intuition that physical things
derive their existence from answered questions, Wheeler's it-from-bit, sits at the root of the whole
computational-universe lineage \cite{wheeler1990}.

The difference is that Vibe Theory takes experience as the substance of the substrate rather than as a
quantity computed from structure. Where integrated information theory measures a property of an organization,
Vibe Theory treats the tone of each cell as a felt state from the start. This is the one step the structural
programs deliberately leave aside, and Vibe Theory is explicit that it is a premise, not a measured result.
Nothing in the experiments measures feeling. What the experiments settle is structure, and the claim that the
structure is felt is a metaphysical commitment placed at the foundation, close in spirit to idealism and to
Wheeler's informational starting point.

\subsection{What Vibe Theory adds overall}

Set against this body of prior work, what is genuinely Vibe Theory's own contribution is narrow and specific.
We state it plainly so as not to oversell it.

\textbf{A forced substrate rather than a random or chosen or assumed one.} Causal sets choose randomness, the
Wolfram program searches a rule space and lets the graph rewrite itself, the triangulations program averages
over an ensemble, and the quantum-cellular-automaton literature assumes a convenient lattice. Vibe Theory
argues the substrate is forced down to a small residue of choice.

Hyperbolic curvature is forced by the demand
for room to grow, the right symmetries, and Lorentz safety at once. Three dimensions are forced by the
existence theorem for compact regular hyperbolic honeycombs and by stable closed orbits. The specific
$\{5,3,4\}$ follows from a stated minimality preference, not a theorem, and Vibe Theory says so. This is a
strong reduction of arbitrariness, not a derivation from nothing.

\textbf{A ternary tone plus the arrow as a value-direction.} The dynamical content is not a rewrite of the
graph but a ternary tone, $-1$, $0$, $+1$, the minimal signed alphabet with opposition and a rest state,
evolving by one local conserving rule. To this Vibe Theory adds one thing none of the prior programs carry,
the arrow, pain below peace below pleasure with desire leaning toward pleasure. The arrow is the single
deliberate posit, a value rather than a derived quantity, and it is what keeps the universe from relaxing to
dead peace. Discrete-physics programs aim at physics. The arrow is what lets the same substrate aim at life,
mind, and intelligence.

\textbf{One substrate that yields physics and life and mind.} The prior programs are theories of physics. Vibe
Theory's larger claim is that the same crystal and the same rule that grow a relativistic quantum field also
grow living, perceiving, intentional systems, with universality (P44, P141) as the capability floor and the
arrow as the direction. This larger claim is more speculative than the physics results, and it is treated as
such.

\textbf{Experiential monism.} Vibe Theory treats the tone of each cell as a felt state from the start.
Universality settles structure and never feeling, and feeling is handled by treating experience as the
substance of the cells rather than as a product of computation. This is a metaphysical commitment, not a
measured result, and it is the part of Vibe Theory furthest outside the established literature.

\textbf{A falsifiable single-substrate experimental program.} Every physics claim above is a number measured
on the same $\{5,3,4\}$ substrate under the same rule, with stated fits and stated negatives. The stochastic
field is reported as massive and diffusive before the unitary completion resolves it. The Lorentz tests
include genuine boosts, not rotation proxies. The crystal step is flagged as a preference and the arrow as a
choice. Much of this reproduces established results to show they live on this one crystal.

In short, Vibe Theory's quantum-field route is the established quantum-cellular-automaton and
quantum-walk-to-Dirac result, run on a hyperbolic substrate whose Lorentz safety is the established
causal-set lesson recast as curvature, with holography and universality inherited from the holographic-code
and hyperbolic-cellular-automaton programs, and a deterministic-substrate spirit inherited from the
cellular-automaton interpretation of quantum mechanics. What is new is the forced substrate, the ternary tone
with the arrow as a value-direction, the experiential monism, and the single-crystal experimental program
that ties physics to life and mind. We claim no more than that.
